\section{Considerações parciais e próximos passos}
\label{sec:consideracoes}

As Fases 1 e 2 transformaram percepções sobre o SIGAA/UnB em uma especificação de avaliação verificável. A Fase 1 delimitou o propósito, as partes interessadas e as decisões apoiadas, adaptou o modelo da ISO/IEC 25010:2023 a uma avaliação caixa-preta e selecionou quatro características pela matriz Impacto~$\times$~Risco. A Fase 2 as desdobrou em quatro objetivos, 14 questões com hipóteses e 20 métricas, com níveis de pontuação e critérios de julgamento fixados antes da coleta, para que o resultado não possa ser ajustado aos dados.

As limitações conhecidas são a amostra de conveniência, concentrada na FCTE; a medição de desempenho apenas no cliente, sujeita a variações de rede e horário; a impossibilidade de avaliar matrícula e trancamento sem alterar dados reais; e a ausência de acesso ao código-fonte.

Na Entrega 2, a Fase 3 produzirá o plano de avaliação: roteiro do moderador e termo de consentimento, publicação do questionário, protocolo das medições e inspeções, ambiente de execução e cronograma. Na Fase 4, as medidas serão obtidas, comparadas aos níveis e julgadas, as hipóteses classificadas e as relações entre as características discutidas. A tarefa e a característica de pior resultado orientarão a ação de melhoria, como um protótipo de alta fidelidade para o fluxo pior avaliado. Na gestão, os marcos e as \textit{issues} serão criados para medir o \textit{burndown}, e a planilha de esforço alimentará os indicadores de recursos e a tabela de contribuição.
